\scanpage{298}
\mspage{39}
l'application $\alpha$ associant à toute $u_i\in\mathrm{Tr}(\unsure{\tau|t},f)$
l'unique intervalle critique qui ``contienne'' $u_i$
(i.e. tel que $u$ soit une flèche dudit intervalle,
ou encore, \unsure{fasse} partie de la ``\unsure{chaîne|chaine} critique''
\unsure{associée}\ldots).

On a mieux encore :

\uline{Proposition 8} Soient $\tau,\tau'\in\mathrm{Ty}$, \deleted{Alors}\inserted{$\tau'$ \uline{réduit}.} Pour \deleted{tout}
morphisme croissant d'int. orientés
\[
 f\colon I_\tau\unsure{\mathbin{-}|\longrightarrow}I_{\tau'}
\]
(on va \deleted{écrire} dire simplement que $f$ est une flèche
de $\underline{\mathrm{Ty}}$, et écrire $f\in\operatorname{Hom}_{\underline{\mathrm{Ty}}}(I_\tau,I_{\tau'})$),
associons lui la partie $\mathrm{Tr}(\tau,f)$ des
flèches de transition de $I_\tau$ non dégénérées par $f$.
Ces flèches forment un ens. \unsure{de} totalement ordonné,
affecté des signes $+$ ou $-$ (caractère direct ou
rétrograde des flèches), et $f$ induit une \inserted{(cor. 5)}
bijection d'ens. tot. ordonnés
\[
\underset{\substack{\cap\\\mathrm{Tr}(\tau)}}{E_f}=\mathrm{Tr}(\tau,f)\simeq\mathrm{Tr}(\tau'),
\]
et $=$ un isom. d'ens. totalement ordonnés
``affectés de signes''. Ainsi, à tout $f$ on associe
une partie $E_f$ de $\mathrm{Tr}(\tau)$, et telle que
son ``type'' soit égal à $\tau'$. Désignons par
\deleted{$\unsure{\mathscr{P}|\mathcal{P}}_{\tau'}(\unsure{\tau})$ la partie de $\mathscr{P}(I_\tau)$}
\deleted{$\mathscr{P}_{\tau'}(\mathrm{Tr}(\tau))$ la partie de}
\[
\mathscr{P}_{\tau'}(\mathrm{Tr}(\tau))\subset\mathscr{P}(\mathrm{Tr}(\tau))
\]
l'ensemble des parties $E$ de $\mathrm{Tr}(\tau)$ dont le type soit $\tau'$.
% archive sheet number: 297
